\documentclass[10pt,letterpaper]{article}
\usepackage{cornell-notes}

\title{Chapter 23: Permanent Function}
\author{}
\date{}

\setDocTitle{Chapter 23: Permanent Function}
\setDocSubtitle{Combinatorial Algorithms}
\setDocVersion{v1.0}
\setDocStatus{Study Companion}
\setDocOwner{Combinatorial Algorithms Cornell Notes}
\setDocAuthor{}
\setDocDate{\today}
\setDocSummary{Cornell-format study and review notes for Chapter 23: Permanent Function.}

\setCornellCollection{Combinatorial Algorithms Cornell Notes}
\setCornellUnitType{Chapter}
\setCornellUnitNumber{23}
\setCornellUnitTitle{Permanent Function}
\setCornellCompanionLabel{Study and review companion}

\hypersetup{pdftitle={Chapter 23: Permanent Function},pdfsubject={Cornell notes},pdfauthor={}}

\begin{document}
\maketitle
\begin{tcolorbox}[colback=Paper,colframe=Ink]{\small\bfseries\color{Teal}CORNELL NOTES}\par{\LARGE\bfseries Chapter 23: The Permanent Function (PERMAN)}\par\medskip\textbf{Topic:} Inclusion-exclusion with Gray-code updates\hfill\textbf{Date:} \today\par\textbf{Study purpose:} Reduce permanent evaluation from factorial to exponential work.\end{tcolorbox}
\begin{tcolorbox}[colback=Teal!5,colframe=Teal,title=Essential question]How do Ryser's inclusion-exclusion formula, complementary subsets, and Gray-code updates yield an $n2^{n-1}$-scale permanent calculation?\end{tcolorbox}
\begin{tcolorbox}[colback=white,colframe=Mist,title=Learning objectives]\begin{itemize}\item Define the permanent and distinguish it from the determinant.\item Interpret $0$-$1$ permanents combinatorially.\item Derive Ryser's subset formula.\item Explain the complementary-subset halving.\item Update row sums along a Gray-code traversal.\item Identify cancellation and overflow risks.\end{itemize}\textbf{Prerequisites:} permutations, inclusion-exclusion, matrices, Gray codes.\end{tcolorbox}
\section*{Cornell notes}
\begin{longtable}{@{}Q!{\color{Mist}\vrule width 1pt}A@{}}\cornellhead
\cue{What is the permanent?}{For an $n\times n$ matrix $A$,
\[
\operatorname{per}(A)=\sum_{\sigma\in S_n}\prod_{i=1}^{n}a_{i,\sigma(i)}.
\]
It resembles the determinant but has no permutation signs. Direct evaluation uses about $n\,n!$ arithmetic operations.}
\cue{What does a $0$-$1$ permanent count?}{It counts permutations selecting only entries equal to one. This includes fixed-point-free permutations for an off-diagonal-one matrix, systems of distinct representatives for an incidence matrix, and legal new rows extending a Latin rectangle.}
\cue{What is Ryser's formula?}{Inclusion-exclusion over missing column images gives
\[
\operatorname{per}(A)=(-1)^n\sum_{S\subseteq[n]}(-1)^{|S|}
\prod_{i=1}^{n}\sum_{j\in S}a_{ij}.
\]
This replaces a sum over $n!$ permutations by $2^n$ subsets.}
\cue{How is the subset count halved?}{Adjoin a constructed column related to row sums. Pair each subset with its complement and choose the added entries so paired products agree. It is then enough to process subsets of $\{1,\ldots,n-1\}$ and double the accumulated result.}
\cue{Why use Gray order?}{For each subset $S$, maintain row quantities $x_i(S)$ formed from a fixed offset plus selected entries. When the next subset differs by column $j$ alone, update every row by $x_i\leftarrow x_i\pm a_{ij}$. This costs $n$ additions rather than recomputing all row sums.}
\cue{What does PERMAN do?}{Initialize the offset row sums, traverse all subsets of the first $n-1$ columns through NEXSUB, use its changed-coordinate output to add or remove one column, form the product of the $n$ current row values, alternate its sign, sum, and apply the final factor $2(-1)^{n-1}$.}
\cue{Cost and storage}{There are $2^{n-1}$ subset states and $O(n)$ update/product work per state, for about $n2^{n-1}$ arithmetic operations. The algorithm uses an $n$-entry Gray bit array and an $n$-entry row-sum array beyond the input matrix.}
\cue{Numeric cautions}{Inclusion-exclusion creates large intermediate terms that cancel. Even an integer matrix can overflow integer storage while the final permanent fits. The source therefore uses double precision. Modern code should consider extended precision, scaling, compensated summation, or exact large integers according to the data.}
\cue{Retrieval practice}{\begin{enumerate}\item What sign is absent from a permanent?\item What does a $0$-$1$ permanent count?\item Why does Ryser use subsets?\item How are complements paired?\item What does the Gray changed index update?\end{enumerate}}
\end{longtable}
\begin{tcolorbox}[colback=Ink!4,colframe=Ink,title=Cornell summary]
The permanent sums one matrix product per permutation and therefore has an expensive direct definition. Ryser's inclusion-exclusion formula converts it to a sum over column subsets. A complementary-subset construction halves the subset family, and a Gray-code traversal makes consecutive subsets differ by one column. PERMAN consequently maintains all row sums with $n$ additions per transition and evaluates roughly $2^{n-1}$ products, reducing work to the $n2^{n-1}$ scale. The method also demonstrates the combinatorial meaning of $0$-$1$ permanents. Its main practical hazard is severe cancellation and large intermediate magnitude, so arithmetic choice matters as much as the combinatorial speedup.
\end{tcolorbox}
\end{document}
